\scanpage{267}
Si \(\deleted{\unsure{C}}\ C\xrightarrow{i}\mathcal M_{N_0}\) est un foncteur exact
\[
\operatorname{Im}\deleted{\unsure{ess}}i\subset\operatorname{Imess}\varphi_{N_0,\infty}
\]
alors on a \(i\simeq\varphi_{N_0,\infty}\tau_{\infty,N_0}i\) d'où
\[
(R^*\varphi_{N_0,\infty})(\tau_{\infty,N_0}\circ i)\simeq R^*i
\]
en particulier
\[
(R^n\varphi_{N_0,\infty})(\tau_{\infty,N_0}i(M))=0\quad\text{si }n>0,\ M\in\operatorname{ob}C.
\]
On est intéressé aux mêmes cas particuliers qu'avant.
\figure{\begin{tikzpicture}\draw (0,0)--(1.4,0);\end{tikzpicture}}
Soit \(K^\bullet\) un complexe de \(C\)

\(L^\bullet\) un complexe de \(\mathcal M_{\deleted{\unsure{\infty}}}\)

et supposons donné un hom dans la catégorie dérivée \(D(\mathcal M_{N_0})\)
\[
i(K^\bullet)\xrightarrow{\alpha}L^\bullet
\]
tel que
\[
\tau_{\infty,N_0}i(K)\xrightarrow[\sim]{\tau_\infty\alpha}\tau_{\infty,N_0}L^\bullet
\]
soit un isom dans \(D(\mathcal M_\infty),\) i.e.

tel que \(\forall n,\)
\diagram{\begin{array}{ccc}
H^n(\alpha):&H^n(i(K))&\longrightarrow H^n(L^\bullet)\\
&\rotatebox{90}{\(\simeq\)}&\\
&i(H^n(K))&
\end{array}}
soit un \(\infty\)-isom i.e. noyau et conoyau sont négligeables. On en conclut
\margin{266}
